\chapter{Byte recognition, binary streams and CFGs}\label{ch:10}
\section*{The working problem}
Select a recognizer according to the shape of the input, not merely the familiarity of its name.

StaticRegex recognizes bytes with a compiled finite-state program. BinaryStream handles explicit byte order. CFG consumes a token vector using Earley-style chart parsing, including left recursion and epsilon productions.

\section*{Protocol selection and chart parsing}
BinaryStream.next recognizes a regex token at the current cursor. It rejects an empty token because a token consumer must advance. Raw reads return views, whereas write may change backing storage. Document those lifetimes in a binary-format wrapper. Serializer and Deserializer are protocols around a user backend's encode and decode operations; they do not choose JSON, CBOR or an on-disk schema for the application.

\section*{In this chapter}
\begin{enumerate}
\item Static byte regexes
\item Explicit endianness
\item Bounded binary reads
\item Token-level context-free parsing
\item Epsilon productions
\end{enumerate}
Read the workshops in order, then use the chapter review to connect the small examples to an application boundary. Each example is independently runnable; it does not depend on variables introduced in an earlier listing.

\clearpage
\workshop{10.1}{Static byte regexes}{C++20}
A static regex compiles its pattern at host compile time and recognizes runtime bytes. The matches operation requires complete consumption. Byte semantics make it suitable for protocols and ASCII-oriented token classes, but do not supply Unicode character properties.

\lstinputlisting[language=C++,caption={Example ex10\_01: C++}]{examples/cpp/ex10_01.inc}
\section*{Execution, step by step}
A compiled pattern is reused for both candidates. The suffix x causes whole-input recognition to fail.

\section*{Read the implementation}
The relevant source entry is StaticRegex in vendor/DSLtk.hpp, around line 10457. Follow the public call into its result construction or dispatch branch. Identify the input ownership, the success condition and the exact failure path before changing this example.

The assertions are the acceptance contract for this workshop. Run the complete program, then change one input at a time. When a change fails, first identify whether the failure belongs to the native build, structural recognition, language checking, execution or the domain result.

\clearpage
\workshop{10.1}{Static byte regexes}{D}
Byte recognition and fixed-width fields must remain independent of native word size. D uint and ushort map to the explicit widths used here; choose byte order at each format boundary rather than inheriting machine endianness.

\lstinputlisting[language=D,caption={Example ex10\_01: D}]{examples/d/ex10_01.inc}
\section*{Failure and design boundary}
Search, prefix and matches have different consumption policies. Use matches for validation and prefix only when maintaining a cursor yourself.

\section*{Try it yourself}
Test an empty candidate and compare plus with star. Explain how zero-length matches affect token loops.

\section*{Review the result}
Keep the original accepted fixture and add the new boundary fixture beside it. Compare observable values and diagnostic categories across the two hosts. Do not label the D run as passed until a D compiler has executed it; the supplied edition records source review separately from runtime evidence.

\clearpage
\workshop{10.2}{Explicit endianness}{C++20}
BinaryStream reads unsigned integers with a selected byte order. This avoids depending on the host CPU layout or pointer alignment. A file format should state field widths and byte order independently of the implementation language.

\lstinputlisting[language=C++,caption={Example ex10\_02: C++}]{examples/cpp/ex10_02.inc}
\section*{Execution, step by step}
The two emitted bytes are 12 and 34 in hexadecimal. Reading with the same order reconstructs the original value.

\section*{Read the implementation}
The relevant source entry is BinaryStream in vendor/DSLtk.hpp, around line 10543. Follow the public call into its result construction or dispatch branch. Identify the input ownership, the success condition and the exact failure path before changing this example.

The assertions are the acceptance contract for this workshop. Run the complete program, then change one input at a time. When a change fails, first identify whether the failure belongs to the native build, structural recognition, language checking, execution or the domain result.

\clearpage
\workshop{10.2}{Explicit endianness}{D}
Byte recognition and fixed-width fields must remain independent of native word size. D uint and ushort map to the explicit widths used here; choose byte order at each format boundary rather than inheriting machine endianness.

\lstinputlisting[language=D,caption={Example ex10\_02: D}]{examples/d/ex10_02.inc}
\section*{Failure and design boundary}
Unsigned width is a format decision. Reading a platform-sized integer can silently change a protocol across architectures.

\section*{Try it yourself}
Read the same bytes as little endian and verify the deliberately different result.

\section*{Review the result}
Keep the original accepted fixture and add the new boundary fixture beside it. Compare observable values and diagnostic categories across the two hosts. Do not label the D run as passed until a D compiler has executed it; the supplied edition records source review separately from runtime evidence.

\clearpage
\workshop{10.3}{Bounded binary reads}{C++20}
A read advances the cursor only when enough bytes are available. Treat truncation as malformed input and retain the offset for a useful error. Borrowed views into stream storage must not outlive mutations that can move that storage.

\lstinputlisting[language=C++,caption={Example ex10\_03: C++}]{examples/cpp/ex10_03.inc}
\section*{Execution, step by step}
Four bytes are consumed and the cursor reaches the buffer end. No integer interpretation is needed for this raw read.

\section*{Read the implementation}
The relevant source entry is BinaryStream in vendor/DSLtk.hpp, around line 10543. Follow the public call into its result construction or dispatch branch. Identify the input ownership, the success condition and the exact failure path before changing this example.

The assertions are the acceptance contract for this workshop. Run the complete program, then change one input at a time. When a change fails, first identify whether the failure belongs to the native build, structural recognition, language checking, execution or the domain result.

\clearpage
\workshop{10.3}{Bounded binary reads}{D}
Byte recognition and fixed-width fields must remain independent of native word size. D uint and ushort map to the explicit widths used here; choose byte order at each format boundary rather than inheriting machine endianness.

\lstinputlisting[language=D,caption={Example ex10\_03: D}]{examples/d/ex10_03.inc}
\section*{Failure and design boundary}
Do not retain a C++ span across a write that can reallocate the vector. Copy the bytes if independent ownership is required.

\section*{Try it yourself}
Attempt a five-byte read from a four-byte stream and translate the exception to a format diagnostic.

\section*{Review the result}
Keep the original accepted fixture and add the new boundary fixture beside it. Compare observable values and diagnostic categories across the two hosts. Do not label the D run as passed until a D compiler has executed it; the supplied edition records source review separately from runtime evidence.

\clearpage
\workshop{10.4}{Token-level context-free parsing}{C++20}
CFG terminals are complete token strings. The lexer decides how source characters become tokens. This separation lets a grammar express recursive structure without embedding numeric scanning and whitespace handling into every production.

\lstinputlisting[language=C++,caption={Example ex10\_04: C++}]{examples/cpp/ex10_04.inc}
\section*{Execution, step by step}
The terminal name matches one token, not four character tokens. A different token string fails at the initial chart position.

\section*{Read the implementation}
The relevant source entry is CFG in vendor/DSLtk.hpp, around line 10671. Follow the public call into its result construction or dispatch branch. Identify the input ownership, the success condition and the exact failure path before changing this example.

The assertions are the acceptance contract for this workshop. Run the complete program, then change one input at a time. When a change fails, first identify whether the failure belongs to the native build, structural recognition, language checking, execution or the domain result.

\clearpage
\workshop{10.4}{Token-level context-free parsing}{D}
D uses module imports and exclamation-mark template instantiation where C++ uses includes, namespaces and angle brackets. The runnable file includes the required imports; the listing shows only its main body.

\lstinputlisting[language=D,caption={Example ex10\_04: D}]{examples/d/ex10_04.inc}
\section*{Failure and design boundary}
An ambiguous CFG returns a first discovered derivation here, not a complete parse forest. Do not claim ambiguity analysis from one returned tree.

\section*{Try it yourself}
Add a production for parenthesized S and parse nested token lists.

\section*{Review the result}
Keep the original accepted fixture and add the new boundary fixture beside it. Compare observable values and diagnostic categories across the two hosts. Do not label the D run as passed until a D compiler has executed it; the supplied edition records source review separately from runtime evidence.

\clearpage
\workshop{10.5}{Epsilon productions}{C++20}
An empty right-hand side denotes epsilon. This allows recursive lists and optional grammar structure, but epsilon should be intentional because it changes which documents the grammar accepts. CFG validation checks names; it does not decide whether optionality matches product requirements.

\lstinputlisting[language=C++,caption={Example ex10\_05: C++}]{examples/cpp/ex10_05.inc}
\section*{Execution, step by step}
The start production completes without consuming a token. A nonempty token stream remains unconsumed and is rejected.

\section*{Read the implementation}
The relevant source entry is CFG in vendor/DSLtk.hpp, around line 10671. Follow the public call into its result construction or dispatch branch. Identify the input ownership, the success condition and the exact failure path before changing this example.

The assertions are the acceptance contract for this workshop. Run the complete program, then change one input at a time. When a change fails, first identify whether the failure belongs to the native build, structural recognition, language checking, execution or the domain result.

\clearpage
\workshop{10.5}{Epsilon productions}{D}
D uses module imports and exclamation-mark template instantiation where C++ uses includes, namespaces and angle brackets. The runnable file includes the required imports; the listing shows only its main body.

\lstinputlisting[language=D,caption={Example ex10\_05: D}]{examples/d/ex10_05.inc}
\section*{Failure and design boundary}
Do not transfer the zero-width-token prohibition from PEG scanning to CFG epsilon productions. The algorithms have different progress mechanisms.

\section*{Try it yourself}
Add a left-recursive list production and compare empty, one-token and three-token cases.

\section*{Review the result}
Keep the original accepted fixture and add the new boundary fixture beside it. Compare observable values and diagnostic categories across the two hosts. Do not label the D run as passed until a D compiler has executed it; the supplied edition records source review separately from runtime evidence.

\clearpage
\section*{Chapter review: Protocol selection and chart parsing}
CFG terminals are lexer-produced token strings. Earley chart items track production, dot and origin, with an attached first-discovered tree. The implementation deduplicates these item coordinates, so it does not retain every ambiguous derivation. Use grammar validation before accepting plugin rules. For a deterministic programming language, specify operator precedence and associativity through productions rather than relying on whichever derivation happens to be discovered first.

\section*{Questions to answer in your implementation}
\begin{enumerate}
\item Test an empty candidate and compare plus with star. Explain how zero-length matches affect token loops.
\item Read the same bytes as little endian and verify the deliberately different result.
\item Attempt a five-byte read from a four-byte stream and translate the exception to a format diagnostic.
\item Add a production for parenthesized S and parse nested token lists.
\item Add a left-recursive list production and compare empty, one-token and three-token cases.
\end{enumerate}
\section*{A practical completion test}
Choose one workshop and replace its fixed input with an application boundary: a command-line argument, a file, a configuration field or a typed callback. Write down the accepted input shape and the error outcome. Keep parsing and validation separate from publishing state. Preserve the smallest failing input when a regression is found.

